%!TEX program = uplatex
\RequirePackage{plautopatch}
\documentclass[uplatex,dvipdfmx]{jlreq}
\usepackage{amsmath,amssymb,amscd,amsthm}

\pagestyle{empty}

\begin{document}

%======================================================================
% 講演１：今吉洋一「リーマン面の退化と再生入門」
%======================================================================
\begin{center}
{\Large\bfseries リーマン面の退化と再生入門}

\hspace*{\fill}{\large\bfseries 今吉 洋一（大阪市立大・理）}
\end{center}

\bigskip

種数 $1$ の閉リーマン面（$=$ トーラス $=$ 楕円曲線）を
具体例にして、今回のテーマである「リーマン面の退化と再生」で
取り上げられる一つの中心的な問題と基本的なことがらとについて
説明する。複素多様体、複素解析、代数幾何、トポロジーなど
様々な分野が交差するところである。

リーマン面の正則族 $\{S_t\}_{t\in R}$ とは、
複素パラメータ $t$ に関して ``正則に動く''
リーマン面の族である。ここで留意することは、各リーマン面 $S_t$ の
位相構造は $t$ に依存しないが、その複素構造は $t$ に関して正則に
変化することである。特に、パラメータ空間 $R$ として、
穴開き円板 $D^* = D \setminus \{0\} =
\{ t \in {\mathbf C} \, | \, 0 < |t| < \delta\}$ の場合を考える。
このとき、素朴な疑問として、
\medskip
\begin{quote}
\bfseries $\bigsqcup_{t \in \Delta} \{t\} \times S_t$
が $2$ 次元複素多様体になるとき、$t=0$ にはどのような
リーマン面 $S_0$ が現れるか？
\end{quote}
\medskip
ということが浮かぶ。具体例での観察から次が分かる。
\begin{itemize}
\item $S_0$ の位相型は一般の $S_t,\ t \neq 0$ とは異なり、
つぶれた（退化した）ものが現れることがある。退化した $S_0$ を
特異ファイバーと呼ぶ。
\item 原点のまわりのモノドロミー $\mathcal{M}$ が重要である。
この $\mathcal{M}$ は、$t$ が原点の廻りを一周するとき、
$S_t$ の標識 $\Sigma_t$（すなわち、その基本群の標準生成元系）の
取り換えを表現する位相的な対象であり、曲面の写像類群の元に対応している。
\end{itemize}
そこで、リーマン面の正則族 $\{S_t\}_{t\in D^*}$ に
対して、次の問題が考えられる。
\begin{itemize}
\item 特異ファイバー $S_0$ を分類せよ。
\item モノドロミー $\mathcal{M}$ を特徴付けよ。
\item $\mathcal{M}$ から正則族 $\{S_t\}_{t\in D}$ を構成せよ
（モノドロミーから正則族の再生）。
\end{itemize}

これらの問題の研究は、小平邦彦先生の楕円曲線の退化の研究を出発点とする。
種数 $2$ の場合は、浪川幸彦氏と上野健爾氏の研究がある。
一般の場合は、松本幸夫氏と J.~M.~Montesinos 氏の共同研究に
よってなされた。これは金曜日の２コマ目のテーマである。
また土曜日には、まず１コマ目では退化ファイバー芽の持つ不変量について、
２コマ目では分裂族とモース化の話題について、
３コマ目ではタイヒミュラー空間の理論を用いた内容について、
そして最後の４コマ目ではレフシェッツ・ファイバー空間について、
それぞれ考察する。

\newpage

%======================================================================
% 講演２：松本幸夫「リーマン面の退化形の位相的分類」
%======================================================================
\begin{center}
{\Large\bfseries リーマン面の退化形の位相的分類}

\hspace*{\fill}{\large\bfseries 松本 幸夫（東大・数理）}
\end{center}

\bigskip

複素単位開円板 $D$ 上に、種数 $g$ のコンパクトリーマン面の
ホロモルフィックな族 $f: M \to D$ があり、原点 $0$ の上にだけ
特異ファイバーがある、しかも極小である（すなわち、特異ファイバーは
$(-1)$-curve を含まない）という状況を考える。これをリーマン面の退化族という。
目標は全ての退化族 $(M, f, D)$ の、ファイバー構造も込めた位相形の分類である。
種数 1 の場合は小平先生による楕円曲面の特異ファイバーの分類がある。
種数 2 の場合には浪川幸彦氏と上野健爾氏による分類がある。ここでは、
種数 $g$ が 2 以上の場合を考える。

もっとも、分類といっても全てを数え上げることを目標にするのではなく
（$g$ が一般では全てを数え上げるのは不可能であろう）、様々な
退化形が現れる位相幾何学的根拠を明らかにすることを目標にする。
そういう立場で「分類」すると、結果は簡単で、
退化族の位相形はそのモノドロミー同相写像の共役類と１対１に対応する、
というものである。ここで、共役類とは写像類群 ${\mathcal{M}}_g$
のなかの共役類のことである。また、モノドロミー
同相写像は「右側ねじれの擬周期写像」として特徴付けられる。
このように、分類する対象はホロモルフィックであるが、様々な退化形
（の複素構造抜きの位相形）は全く
位相幾何学的な現象として現れていることが分かる。
以上の結果の証明の概略を紹介したい。

\newpage

%======================================================================
% 講演３：足利正「退化不変量の様々な顔」
%======================================================================
\begin{center}
{\Large\bfseries 退化不変量の様々な顔}

\hspace*{\fill}{\large\bfseries 足利 正（東北学院大・工）}
\end{center}

\bigskip

退化を持つリーマン面族の全空間となっている
コンパクト複素曲面もしくは実４次元多様体の
大域的不変量の「局所版」としての退化ファイバー芽の不変量は、
歴史的にはまず種数が２の時に、一般型曲面論の立場から
堀川頴二氏（1976）、保形形式的観点から上野健爾氏（1987）、
位相幾何的立場より松本幸夫氏（1996）によりそれぞれ提出された。
しかしこれらは異なる三種の不変量を与えるのではなく、実は本質的には
同一の不変量の互いに相異なる三つの「顔」なのである。

ここ数年、これらの不変量のある意味での高種数への拡張
と思えるものが、次々と提案された。
一は今野一宏氏、二は吉川謙一氏、三は（遠藤久顕氏、森藤孝之氏
のものの拡張と思える）古田幹雄氏によるものである。
これらも同じ不変量の相異なる顔なのかもしれないが、
その事はまだ証明されておらず、またその性質の探求は今後に託されている。

ここでは、これらを網羅的に紹介するのは筆者の能力の
及ぶ所ではないので、ある観点からこの方面の若干の説明を試みたい。

\newpage

%======================================================================
% 講演４：高村茂「分裂族の織りなす宇宙」
%======================================================================
\begin{center}
{\Large\bfseries 分裂族の織りなす宇宙}

\hspace*{\fill}{\large\bfseries 高村 茂（京大・理）}
\end{center}

\bigskip

私見であるが、リーマン面の「退化」という用語は後ろ向きの感じがする。
音はそのままで「大化」というのはどうだろう。
ファイバーの形が大きく変化することをピッタリ表わしている感じがするが。

冗談はさておき、
リーマン面の退化を摂動すると、もともと一つだった特異ファイバーが
いくつかの特異ファイバーに分裂することがある。
このような摂動を分裂族と呼んでいる。
まったく分裂族をもたない退化もまれに存在し、原子退化と呼ばれる。

分裂族に関する主要な問題としては（1）原子退化の分類、
（2）与えられた退化のすべての分裂族を記述すること、などが挙げられる。
（1）の分類は種数 5 以下のリーマン面の退化の場合に完成しており、
種数 1、2 はそれぞれ
モイシェゾン氏、堀川氏、種数 3、4、5 は私による。
なお、種数 1、2 のときに使われた２重分岐被覆法は足利、荒川両氏に
より超楕円的リーマン面の退化の場合へと発展した。
これに対し種数 3 以上の
リーマン面は一般に超楕円的ではないが、
私は別の手法で特異ファイバーの形が星型でないときは
種数にかかわらず、（例外的な場合を除いて）退化は必ず分裂することを示した。
さらに、私の開発した「剥がし変形」という分裂族の構成法は、
特異ファイバーが星型の場合にも適用できる。
種数 3、4、5 の場合の原子退化の分類は、これらの結果を用いて得られる。

（2）のどのような分裂がおこりうるか？という問いから
自然に「分裂のモジュライ空間」の構成問題へと導かれる。
「分裂のモジュライ空間」の候補として、
ある複素解析空間の族を構成したが、
そこではベクトル束の対称積をひねったものが活躍する。

将来的には上の複素解析空間の族の「解析的複雑さ」と、
退化のモノドロミーの写像類群の元としての「組み合わせ的複雑さ」
をつなぐ何らかの原理がみつかり、分裂族の織りなす宇宙の姿が
見えてくるのではと期待している。
それこそ「大化」の改新である。

\newpage

%======================================================================
% 講演５：今吉洋一「タイヒミュラー空間から見たリーマン面の退化と再生」
%======================================================================
\begin{center}
{\Large\bfseries タイヒミュラー空間から見たリーマン面の退化と再生}

\hspace*{\fill}{\large\bfseries 今吉 洋一（大阪市立大・理）}
\end{center}

\bigskip

「リーマン面の退化と再生入門」で取り上げた問題を、
タイヒミュラー空間 $T_{g, n}$ の理論を用いて考察したい。
特に、穴開き円板 $D^* = \{ 0 < |t| < \delta \}$ 上の
正則族 $\{S_t\}_{t\in D^*}$ の原点でのモノドロミー $\mathcal{M}$ の
特徴付けを取り上げる。このとき複素解析的な対象で
ある $\{S_t\}_{t\in D^*}$ の性質と、位相的な対象で
ある $\mathcal{M}$ の性質との相互関係に興味がある。

我々の考察にタイヒミュラー空間 $T_{g, n}$ が使われる背景には、
$T_{g, n}$ はモデュライ空間 $M_{g, n}$ の普遍被覆空間に相当し、
その普遍被覆変換群が写像類群 ${\mathrm{Mod}}_{g, n}$
（$=$ タイヒミュラー・モデュラー群）となることがある。
従って、$M_{g, n} = T_{g, n}/ {\mathrm{Mod}}_{g, n}$ と表わせる。
このことから、正則族 $\{S_t\}_{t\in D^*}$ は $D$ から
$T_{g, n}$ への正則写像として表現され、$\mathcal{M}$ も $T_{g, n}$ への
作用（自己双正則写像）としてとらえることができる。我々の考察においては、
クライン群、（実と複素の）双曲幾何も重要な役割を果たす。

\newpage

%======================================================================
% 講演６：松本幸夫「Lefschetz fibration の紹介」
%======================================================================
\begin{center}
{\Large\bfseries Lefschetz fibration の紹介}

\hspace*{\fill}{\large\bfseries 松本 幸夫（東大・数理）}
\end{center}

\bigskip

リーマン面の退化形のうち、もっとも簡単なものは、
そのモノドロミー同相写像が右側ねじれの Dehn twist になっている
ような特異ファイバーであろう。これが Lefschetz 型の特異ファイバーである。
「特異性として唯１個のノードだけを
持つ特異ファイバー」といっても同じことである。そして、向きの
ついた実 4 次元閉多様体 $M$ から向きのついた閉曲面 $B$ への滑らかな写像
$f: M \to B$ が Lefschetz fibration であるとは、それが、
有限個の特異ファイバーを除いて、種数 $g$ の閉曲面を
ファイバーとするファイバー束になっており、特異ファイバーは
Lefschetz 型のみであることである。このようなファイバー空間（というより、
「ペンシル」）は1920年代に S.~Lefschetz が代数曲面の位相的研究の補助手段として
導入したものであるが、最近、それが 4 次元シンプレクティック多様体
に適合した位相的構造であることが証明され（Donaldson、Gompf）、
多くの人々の注目するところとなった。いろいろな人がいろいろな研究を
していて、とてもそれらを紹介する能力はないが、
面白いと思われるいくつかの結果を紹介したい。

\end{document}
